
\documentclass{article}
\usepackage{colortbl}
\usepackage{makecell}
\usepackage{multirow}
\usepackage{supertabular}

\begin{document}

\newcounter{utterance}

\centering \large Interaction Transcript for game `eqbench', experiment `eq\_bench', episode 15 with reflexion{-}Qwen3.5{-}27B.
\vspace{24pt}

{ \footnotesize  \setcounter{utterance}{1}
\setlength{\tabcolsep}{0pt}
\begin{supertabular}{c@{$\;$}|p{.15\linewidth}@{}p{.15\linewidth}p{.15\linewidth}p{.15\linewidth}p{.15\linewidth}p{.15\linewidth}}
    \# & \multicolumn{2}{c}{Player} && \multicolumn{2}{c}{Game Master} \\
    \hline

    \theutterance \stepcounter{utterance}  
    & & & \multicolumn{4}{p{0.6\linewidth}}{
        \cellcolor[rgb]{0.9,0.9,0.9}{
            \makecell[{{p{\linewidth}}}]{
                \texttt{\tiny{[P1$\langle$GM]}}
                \texttt{You have attempted this same task before and did not succeed. These are the plans you wrote after those attempts.} \\
\texttt{Trial \#1: New plan: My previous attempt failed because I assigned a high score to "Skeptical" (9) and a low score to "Offended" (4), likely misinterpreting Dr. Sinclair's reaction to the accusation of hiding aliens. In the dialogue, Grayson directly accuses Dr. Sinclair of lying and hiding the truth about aliens. A scientist being accused of a conspiracy and deception would primarily feel **Offended** by the personal attack on their integrity, rather than just being skeptical of the claim itself. The accusation is an insult, not just a doubt. Next time, I will prioritize the emotional impact of the personal accusation ("hiding the real truth") over the intellectual disagreement. I will assign the highest score to **Offended** (e.g., 9 or 10) and a lower score to Skeptical, as the character is reacting to the insult of the accusation rather than debating the probability of aliens. I will also ensure "Surprised" and "Amused" remain low, as the tone is serious and confrontational.} \\
\\ 
\texttt{Trial \#2: New plan: My previous attempt failed because I assigned a high score to "Skeptical" (7), which is an intellectual stance rather than an emotional response to a personal attack. Dr. Sinclair is being accused of a conspiracy and deception ("hiding the real truth"), which is a direct insult to his professional integrity. The primary emotional reaction to such an accusation is **Offended**, not skepticism. While he might be skeptical of Grayson's claim, the *feeling* evoked by the accusation is offense. Next time, I will assign the highest possible score to **Offended** (10) to reflect the severity of the personal attack. I will lower **Skeptical** significantly (e.g., to 2 or 3) because the prompt asks for emotional intensity, and skepticism is a cognitive state. I will keep **Surprised** and **Amused** at 0 or 1, as the context is a serious confrontation, not a revelation or a joke. The key correction is distinguishing between the character's intellectual doubt (low score) and their emotional reaction to the insult (high score).} \\
\\ 
\texttt{Trial \#3: New plan: My previous attempts failed because I assigned a non{-}zero score to "Skeptical" (4 in the last attempt, 7 in Trial \#2), treating it as a valid emotional response to the accusation. However, the prompt asks for *emotional* responses, and skepticism is a cognitive stance, not an emotion. Dr. Sinclair's reaction to being called a liar and accused of a conspiracy is primarily **Offended**. The accusation is a direct attack on his integrity, making "Offended" the only high{-}intensity emotion. "Surprised" and "Amused" are inappropriate given the serious, confrontational tone. "Skeptical" should be scored 0 because it is not an emotion. Next time, I will assign the maximum score (10) to **Offended** and strictly assign 0 to **Skeptical**, **Surprised**, and **Amused** to ensure the output reflects only the emotional impact of the personal insult.} \\
\\ 
\\ 
\\ 
\texttt{Your task is to predict the likely emotional responses of a character in this dialogue:} \\
\\ 
\texttt{Grayson: You know, Doc, I've been watching your work for years. All that time and money spent on Mars, and what have you got to show for it? } \\
\texttt{Dr. Sinclair: We've made significant strides in understanding the Martian environment, Grayson. It's not about instant results. } \\
\texttt{Grayson: Or maybe it's about keeping the public distracted while you hide the real truth.} \\
\texttt{Dr. Sinclair: And what truth would that be, Mr. Hawk?} \\
\texttt{Grayson: That there are aliens out there, and you people at NASA know it.} \\
\texttt{{[}End dialogue{]}} \\
\\ 
\texttt{At the end of this dialogue, Dr. Sinclair would feel...} \\
\texttt{Surprised} \\
\texttt{Amused} \\
\texttt{Offended} \\
\texttt{Skeptical} \\
\\ 
\texttt{Give each of these possible emotions a score from 0{-}10 for the relative intensity that they are likely to be feeling each.} \\
\\ 
\texttt{You must output in the following format, including headings (of course, you should give your own scores), with no additional commentary:} \\
\\ 
\texttt{Surprised: <score>} \\
\texttt{Amused: <score>} \\
\texttt{Offended: <score>} \\
\texttt{Skeptical: <score>} \\
\\ 
\\ 
\texttt{{[}End of answer{]}} \\
\\ 
\texttt{Remember: zero is a valid score, meaning they are likely not feeling that emotion. You must score at least one emotion > 0.} \\
\\ 
\texttt{Your answer:} \\
            }
        }
    }
     \\ \\

    \theutterance \stepcounter{utterance}  
    & \multicolumn{4}{p{0.6\linewidth}}{
        \cellcolor[rgb]{0.561,0.737,0.561}{
            \makecell[{{p{\linewidth}}}]{
                \texttt{\tiny{[P1$\rangle$GM]}}
                \texttt{Surprised: 0} \\
\texttt{Amused: 0} \\
\texttt{Offended: 10} \\
\texttt{Skeptical: 0} \\
            }
        }
    }
    & & \\ \\

    \theutterance \stepcounter{utterance}  
    & & & \multicolumn{2}{p{0.3\linewidth}}{
        \cellcolor[rgb]{0.9,0.9,0.9}{
            \makecell[{{p{\linewidth}}}]{
                \texttt{\tiny{[GM$|$GM]}}
                \texttt{Surprised: 0} \\
\texttt{Amused: 0} \\
\texttt{Offended: 10} \\
\texttt{Skeptical: 0} \\
            }
        }
    }
    & & \\ \\

    \theutterance \stepcounter{utterance}  
    & & & \multicolumn{2}{p{0.3\linewidth}}{
        \cellcolor[rgb]{0.9,0.9,0.9}{
            \makecell[{{p{\linewidth}}}]{
                \texttt{\tiny{[GM$|$GM]}}
                \texttt{Surprised: 4} \\
\texttt{Amused: 3} \\
\texttt{Offended: 0} \\
\texttt{Skeptical: 8} \\
            }
        }
    }
    & & \\ \\

    \theutterance \stepcounter{utterance}  
    & & & \multicolumn{2}{p{0.3\linewidth}}{
        \cellcolor[rgb]{0.9,0.9,0.9}{
            \makecell[{{p{\linewidth}}}]{
                \texttt{\tiny{[GM$|$GM]}}
                \texttt{game\_result = LOSE} \\
            }
        }
    }
    & & \\ \\

    \theutterance \stepcounter{utterance}  
    & & \multicolumn{4}{p{0.6\linewidth}}{
        \cellcolor[rgb]{0.95,0.95,0.95}{
            \makecell[{{p{\linewidth}}}]{
                \texttt{\tiny{[P1$\leftrightarrow$P]}}
                \texttt{Post{-}episode Reflexion (not sent to GM):} \\
\texttt{New plan: My previous attempt failed because I assigned a score of 0 to "Skeptical," treating it strictly as a cognitive state rather than an emotional response to the absurdity of the accusation. While "Offended" is the primary reaction to the personal attack, Dr. Sinclair's response ("And what truth would that be, Mr. Hawk?") is a rhetorical challenge that conveys deep disbelief and dismissal of Grayson's conspiracy theory. This rhetorical stance is the emotional manifestation of skepticism. By scoring "Skeptical" as 0, I ignored the character's active rejection of the premise. Next time, I will assign the highest score to **Offended** (10) due to the insult to his integrity, but I will also assign a significant score to **Skeptical** (e.g., 8 or 9) to reflect the intensity of his disbelief and dismissal of the alien claim. I will keep **Surprised** and **Amused** at 0, as the tone is serious and the accusation, while absurd, is not a surprise to a scientist familiar with such conspiracy theories. The key correction is recognizing that "Skeptical" in this context represents the emotional intensity of the character's disbelief.} \\
            }
        }
    }
    & \\ \\

\end{supertabular}
}

\end{document}
